\section{切平面为何是全局下界}
\label{sec:09-Convex-Optimization/03-Convex-Functions/02}

一个资源分配问题交给求解器，日志的最后两行是：目标值 \(12438.2\)，下界 \(12431.6\)，相对间隙 \(0.05\%\)。第一行是找到的那个方案，可以逐项核对。第二行不一样：它宣称整个可行域里不存在目标值低于 \(12431.6\) 的方案。可行域是连续的，有无穷多个点，求解器只访问了其中几百个，它凭什么替剩下的所有点作担保？

同一个问题在训练场景里换一副面孔。逻辑回归跑到梯度范数 \(10^{-8}\)，可以收工了；同样的数字出现在一个深层网络上，只说明附近这一小块平坦，谁也不知道换个初始化会不会低一大截。两句话的落差不来自算法，也不来自数据规模，它全部来自本篇要建立的那一条不等式。

先看它在最简单的函数上是什么样子。\(f(x)=x^2\) 在 \(x=2\) 处的切线是 \(T(x)=4+4(x-2)\)。在 \(x=0\) 处 \(T(0)=-4\)，而 \(f(0)=0\)；在 \(x=5\) 处 \(T(5)=16\)，而 \(f(5)=25\)；在 \(x=-1\) 处 \(T(-1)=-8\)，而 \(f(-1)=1\)；在 \(x=10\) 处 \(T(10)=36\)，而 \(f(10)=100\)。跑多远都不越界。换成 \(f(x)=\sin x\)，在 \(x=\pi/2\) 处切线是水平线 \(y=1\)，走到 \(x=3\pi/2\)，函数值是 \(-1\)，切线早已穿进图像里面。

\begin{center}
\begin{tikzpicture}[x=1.3cm,y=0.95cm,font=\small,>=stealth]
  \draw[->,black!55] (-0.6,0) -- (3.35,0) node[below,font=\footnotesize] {\(x\)};
  \draw[blue!65!black,very thick] plot[domain=-0.35:3.05,samples=75] (\x,{0.9*(\x-1.3)*(\x-1.3)+1.0});
  \draw[red!70!black,thick] (0.7,0.028) -- (3.05,2.566);
  \fill[red!70!black] (1.9,1.324) circle (2.2pt);
  \node[font=\footnotesize,red!70!black,anchor=south east] at (1.85,1.38) {\(x\)};
  \draw[<->,black!55] (0.75,0.082) -- (0.75,1.272);
  \node[font=\footnotesize,black!70,anchor=east] at (0.62,0.68) {余量 \(\ge0\)};
  \node[font=\footnotesize,align=center,blue!65!black] at (1.4,3.75) {凸：切线不越过曲线};
  \begin{scope}[shift={(6.0,0)}]
    \draw[->,black!55] (-2.1,0) -- (2.3,0) node[below,font=\footnotesize] {\(x\)};
    \draw[blue!65!black,very thick] plot[domain=-1.9:2.0,samples=90] (\x,{0.3*\x*\x*\x*\x-\x*\x-0.35*\x+2.0});
    \draw[red!70!black,thick] (-2.0,1.602) -- (2.1,1.602);
    \fill[red!70!black] (-1.19,1.602) circle (2.2pt);
    \node[font=\footnotesize,red!70!black,anchor=south] at (-1.19,1.70) {\(x\)};
    \fill[blue!65!black] (1.39,0.701) circle (2.0pt);
    \draw[<->,black!55] (1.39,0.701) -- (1.39,1.602);
    \node[font=\footnotesize,black!70,anchor=west] at (0.45,0.42) {曲线低于切线};
    \node[font=\footnotesize,align=center,blue!65!black] at (0,3.75) {非凸：切线穿进曲线};
  \end{scope}
\end{tikzpicture}

\emph{左边这条直线是在一个点上算出来的，却对每一个 \(y\) 都成立；右边同样一条切线，走远一点就失去了担保。}
\end{center}

图里左边那条红线，是这一整个部分的枢纽。它由一个点的函数值和梯度完全确定——两样都是局部量、都能当场算出来——而它给出的断言覆盖整个定义域。凸优化后面所有的威力，对偶、停止准则、``梯度为零即最优''，追到底都是这一条。

\subsection{一阶下界与弦定义是同一件事}

先把它写清楚。设 \(f\) 在开凸集上可微，则对定义域中所有 \(x,y\)，

\[
f(y)\ge f(x)+\nabla f(x)^{\trans}(y-x).
\]

右边是 \(f\) 在 \(x\) 处的一阶 Taylor 展开。对一般光滑函数，这个表达式只是 \(y\to x\) 时的近似，误差随距离增长；对凸函数，它在整个定义域上是一个有效的下界。

证明只有两步。把 \(f\) 限制在连接两点的线段上，令 \(g(t)=f(x+t(y-x))\)，\(t\in[0,1]\)。凸性遗传给 \(g\)，于是 \(g(t)\le(1-t)g(0)+tg(1)\)，移项得

\[
g(1)\ge g(0)+\frac{g(t)-g(0)}{t}.
\]

第二步令 \(t\downarrow0\)。右端的差商收敛到 \(g'(0)=\nabla f(x)^{\trans}(y-x)\)，代入 \(g(0)=f(x)\)、\(g(1)=f(y)\) 即得。

反向也只有两步。设每一点的切平面都是全局下界，取线段上的点 \(z=\theta x+(1-\theta)y\)，分别以 \(x\) 和 \(y\) 为参照写两条下界，第一条乘 \(\theta\)、第二条乘 \(1-\theta\) 再相加。由于 \(\theta(x-z)+(1-\theta)(y-z)=0\)，梯度项整体消失，剩下 \(\theta f(x)+(1-\theta)f(y)\ge f(z)\)，正是弦不等式。

条件对一遍。可微性只在 \(x\) 处用到，用的还只是沿 \(y-x\) 方向的方向导数；定义域取开集是为了让 \(t\downarrow0\) 时 \(g\) 有定义，边界点要单独处理。中间那个不等式对每个 \(t\in(0,1]\) 都成立，取极限不需要额外的连续性假设——凸函数在开域内部自动连续。可微性失效时结论并不作废，只是梯度要换成后面的次梯度。

\subsection{梯度为零就是一张全局证书}

把 \(\nabla f(x^\star)=0\) 代进一阶下界，右边的线性项消失：

\[
f(y)\ge f(x^\star)\qquad\text{对所有 }y.
\]

\(x^\star\) 是全局最优点。不用检查 Hessian，不用比较别的候选点，不用担心还有一个没搜到的盆地。一个在单点上算出来的量，直接判决了整个空间。

非凸情形下同一个数字只标注一个驻点，可能是极小、极大或鞍点，还可能是一片高原。差别不在于凸函数的梯度携带了更多信息——两边算的都是同一个偏导数向量——而在于凸性把``此处平坦''和``全局没有更低点''这两句话接上了。少了那条下界，第一句推不出第二句。

这也回答了开头日志里的第二行。求解器报的下界并不是从可行点里挑出来的，它来自一族仿射函数：每一个都全局地压在目标函数下面，因此每一个的最小值都是原问题最优值的合法下界。挑一个最紧的仿射下界，就得到对偶问题，这条路线的完整展开见\hyperref[sec:09-Convex-Optimization/05-Duality/01]{``Lagrangian 怎样制造全局下界''}。有了下界，``还能不能更好''就成了一个可以量化的问题：算出的目标值减去下界，就是当前解距离最优的上限，这个差额被称为对偶间隙，它是凸求解器能给出停止准则的原因，见\hyperref[sec:09-Convex-Optimization/05-Duality/02]{``对偶间隙何时闭合''}。

拿它跟非凸训练对照就更清楚了。深度网络里没有这样的证书，只能拿梯度范数、损失曲线的斜率、验证集指标凑合着判断该不该停；\hyperref[sec:09-Convex-Optimization/08-Nonconvex-and-Stochastic-Optimization/01]{``下降引理在非凸地形中只保证接近驻点''}给出的最强结论也只是``离某个驻点足够近''。凸性买到的东西，一句话概括就是：局部检查可以结案。

\subsection{约束把零梯度换成没有可行的下降方向}

在凸集 \(C\) 上最小化可微凸函数时，最优点未必落在梯度为零的地方，条件要改写成

\[
\nabla f(x^\star)^{\trans}(y-x^\star)\ge0
\qquad\text{对所有 }y\in C.
\]

它说的是：从 \(x^\star\) 指向任何一个可行点的方向，一阶变化率都不为负。所有允许的移动要么平坦要么上坡，于是没有可行的改进方向。

充分性还是那条下界：对任意可行 \(y\)，\(f(y)\ge f(x^\star)+\nabla f(x^\star)^{\trans}(y-x^\star)\ge f(x^\star)\)。无约束是特例——\(C=\R^n\) 时可以朝正负坐标方向各走一步，两个不等式合起来逼出梯度的每个分量为零。在边界上就不同了，梯度可以不为零，只要它指向的下降方向被可行域挡住。这个几何图像在\hyperref[sec:09-Convex-Optimization/05-Duality/03]{``KKT 把最优性写成一组方程''}里会被翻译成乘子的语言：挡住下降方向的那面墙，收多少``过路费''，就是那个约束的乘子。

\subsection{二次函数把余量算得一清二楚}

一阶下界留了多少余地，在二次函数上可以直接读出来。取

\[
f(x)=\tfrac12x^{\trans}Qx+b^{\trans}x+c,
\qquad Q\succeq0,
\]

其梯度为 \(\nabla f(x)=Qx+b\)。把差值展开：

\[
\begin{aligned}
f(y)-f(x)-\nabla f(x)^{\trans}(y-x)
&=\tfrac12y^{\trans}Qy+b^{\trans}y-\tfrac12x^{\trans}Qx-b^{\trans}x-(Qx+b)^{\trans}(y-x)\\
&=\tfrac12(y-x)^{\trans}Q(y-x).
\end{aligned}
\]

余量恰好是以 \(Q\) 为度量的两点距离的平方的一半。\(Q\succeq0\) 保证它非负，半正定与一阶下界在这里是同一个条件的两种写法。

余量的大小随即决定了别的东西。\(Q\succ0\) 时余量在任何方向上都严格为正，\(Qx^\star+b=0\) 有唯一解且是全局最优。\(Q\) 只是半正定时，零空间方向上余量为零，目标沿那些方向完全平坦：若 \(b\) 在零空间上的分量为零，最优解是一整个仿射子空间；若不为零，目标无下界。岭回归往 \(Q\) 上加 \(\lambda I\) 做的正是消灭这片平坦区，代价与收益见\hyperref[sec:09-Convex-Optimization/07-Applications/01]{``光滑经验风险与 \(L_2\) 正则化''}。

\subsection{不可微点上支持平面照样存在}

\(f(x)=\abs{x}\) 在原点没有切线，左边割线斜率是 \(-1\)，右边是 \(+1\)。但下界这件事并没有随之消失：任何一条过原点、斜率在 \([-1,1]\) 之间的直线，都整条位于 \(\abs{x}\) 的图像下方。可用的支撑线不是没有了，是多了。

把这个观察当作定义。若向量 \(g\) 满足

\[
f(y)\ge f(x)+g^{\trans}(y-x)
\qquad\text{对所有 }y,
\]

则称 \(g\) 是 \(f\) 在 \(x\) 处的\textbf{次梯度}，全体次梯度构成的集合称为\textbf{次微分}，记作 \(\partial f(x)\)。注意这个定义里没有出现极限，它直接把上面那条不等式抬举成了定义本身。对绝对值函数，

\[
\partial f(x)=
\begin{cases}
\set{1},&x>0,\\[2pt]
[-1,1],&x=0,\\[2pt]
\set{-1},&x<0.
\end{cases}
\]

可微点处次微分退化成单点集 \(\set{\nabla f(x)}\)，尖点处它是一个集合。最优性条件跟着推广为 \(0\in\partial f(x^\star)\)：对 \(\abs{x}\) 而言 \(0\in[-1,1]=\partial f(0)\)，所以 \(x^\star=0\) 是全局最优，尽管这里导数不存在。\(\ell_1\) 正则化的整套稀疏性论证就靠这个集合的宽度撑着，见\hyperref[sec:09-Convex-Optimization/07-Applications/02]{``\(L_1\) 正则为何产生稀疏''}；算法层面的用法见\hyperref[sec:09-Convex-Optimization/06-Optimization-Algorithms/03]{``不可微时次梯度为何仍能前进''}。

有一点值得提醒。次梯度不是``导数不存在时随便凑一个近似梯度''，它必须逐点满足那条全局不等式。非凸场景里也有各种广义梯度的定义，形式上相似，但它们不继承全局下界的含义，因而也不提供证书。自动微分在 \(\relu\) 的折点处返回一个数，那个数在凸情形下确实是一个合法次梯度，在非凸网络里就只是一个约定。

\subsection{三个下界是同一条支撑线}

把几张图叠起来。上境图是凸集，函数图像上的点是它的边界点，凸集在边界点处存在支持超平面——这是上一单元分离定理的直接推论。支持超平面的斜率就是次梯度，超平面本身就是那条不越过图像的仿射函数。于是凸函数的一阶下界、凸集的支持超平面、对偶问题给出的全局下界，三者找的是同一样东西：一条不穿进凸对象内部的线性支撑。切平面支撑的是函数图像，分离超平面支撑的是可行域，Lagrange 乘子是把两种支撑合成一条的组合系数。

这个统一说明了为什么仿射下界值得单独立一节。假如你能为所有可行点造出一个下界，并且在手上这个候选解处下界恰好贴住目标值，那么解和最优性证明就同时到手了，不需要继续搜索，也不需要留一份怀疑。剩下的工作是把这件事做得系统、做得可计算，那正是对偶理论要处理的。

一阶下界保证了``到了平坦处就不会更差''。它没说最优解唯一，也没说走到那里要花多久——余量项 \(\tfrac12(y-x)^{\trans}Q(y-x)\) 只被要求非负，可以要多小有多小。下一篇把这个余量本身当作研究对象：它严格为正带来什么，它被一个正数从下面顶住又带来什么。
